\section{The density clause and a precise stabilizer converse}
\label{sec:plethysm-density-converse}

\paragraph{The exact narrow question remains open.}
Write $\ell(a)=1+\lfloor\log_2a\rfloor$ for each positive
partition part and add these lengths over the three partitions;
zero parts are omitted.  Let $\mathcal D_N$ be the ordinary
plethysm triples of total length exactly $N$ with nonzero initial
coefficient, and let $\mathcal E_N$ be those whose literal
saturation index is positive.  The conditional reading of
GCT VI, Hypothesis 1.6.5(b), printed p.~25, is
\[
 \text{for every fixed }c>0,\qquad
          |\mathcal E_N|/|\mathcal D_N|\le N^{-c}
                  \quad\text{for all sufficiently large }N.
 \tag{PNEW-density-target}
\]
The new positive initial example proves that this exceptional
set is nonempty; it does not settle its asymptotic density.
The full inherited proof P-F48, reproduced below and rechecked
here, bounds all the cubic rescalings and their inner and outer
determinant twists by an exponentially small proportion.
No map with logarithmic raw-length overhead from arbitrary
ordinary plethysm inputs to such exceptions has been constructed.
Consequently this work neither proves nor refutes
(PNEW-density-target).  This is an explicit unresolved part of
the user request, not an inference from the refutation of SH(a).
PNEW-009 below does prove this exact-shell conclusion after
fixing the inner representation and rank, with constants
depending on that fixed embedding.  Summing over changing
embeddings remains a separate uniformity problem.

\begin{proposition}[PNEW-008: no input-growth converse for the intrinsic graded quotient]
For each integer $b\ge0$ take the ordinary plethysm input
\[
 \mu=(3),\quad
 \lambda_b=(35+b,b^{679}),\quad
 \pi_b=((7+136b)^{15}),
 \tag{PNEW-twisted-positive}
\]
with zero parts omitted.  These inputs have unbounded raw
partition length, initial coefficient one, and strict saturation
index at least seven.  All their canonical graded multiplicity
algebras are isomorphic, compatibly with the grading.
In particular the orders of all stabilizers for the intrinsic
grading action on their nonzero affine quotients are bounded by
one fixed finite constant, independent of $b$.

Thus a positive saturation index, even with positive initial
coefficient, does not force grading stabilizer orders to grow
with the encoded input.  This is a statement about the fixed
effective graded quotient, not about every possible affine
presentation and its raw general-linear determinant character.
\end{proposition}
\begin{proof}
Let $V=\mathbb C^{15}$ and $W=\operatorname{Sym}^3V$.
Then $\dim W=680$ and $\det W=(\det V)^{136}$.
Indeed each variable has equal total exponent in the monomial
basis, and the sum of those exponents is $3\cdot680$.
For every $s\ge0$ the determinant-twist identity gives
\[
 S^{s\lambda_b}W
     =\operatorname{Sym}^{35s}W\otimes(\det W)^{bs}
     =\operatorname{Sym}^{35s}W\otimes(\det V)^{136bs}.
\]
Cancelling the same character from the target
$(\det V)^{(7+136b)s}$ identifies the multiplicity space
with that of $(\det V)^{7s}$ in $\operatorname{Sym}^{35s}W$.
These identifications respect Cartan multiplication: the chosen
one-dimensional character vectors multiply by tensor product
and their common factors cancel degree by degree.  The whole
graded algebra is therefore the single ring with Hilbert
function $h(s)=f_{15}(7s)$ of PNEW-007.  Its initial coefficient
and index bounds follow from that proposition.  The lengths
are unbounded, since the positive parts tend to infinity with $b$.

For the stabilizer assertion, choose homogeneous algebra
generators $u_1,\ldots,u_t$ of this fixed finitely generated
graded ring, and put $D=\max_i\deg u_i$.  At any nonzero
point $y$ of its affine cone, some positive-degree generator
does not vanish: otherwise every positive-degree function
would vanish and $y$ would be the graded origin.  A scalar
$z\in\mathbb C^*$ fixes $y$ precisely when
$z^{\deg u_i}=1$ for all the nonzero values $u_i(y)$.
Its stabilizer is therefore
\[
                  \boldsymbol\mu_{\gcd\{\deg u_i:u_i(y)\ne0\}},
\]
and has order at most $D$.  The bound is uniform over all
points and all $b$, because the ring and its grading are the
same.  It applies in particular to every stable point for
this grading action on the nonzero locus.
\end{proof}

\paragraph{The simpler fixed Kahle--Micha\l ek hole.}
The same construction applied to the fixed ray
$\phi(s)=a_{(4s),(3)}^{(7s,5s)}$ gives
\[
       \lambda_b=(4+b,b,b,b),\qquad
       \pi_b=(7+6b,5+6b),\qquad\mu=(3).
\]
Its stretching algebra is unchanged, since
$\det(\operatorname{Sym}^3\mathbb C^2)=(\det\mathbb C^2)^6$.
The checked inherited identification PLE-F8 realizes it as
$\mathbb C[U,V,Z]/(Z^2-UV^2-4U^4)$ with degrees $(2,3,4)$.
Hence all its nonzero grading stabilizers have order at most
4, while $p(\phi)=s(\phi)=2$ and the input lengths are unbounded.
This simpler example has $\phi(1)=0$, so it is excluded from
the conditioned population in (PNEW-density-target).  PNEW-008
removes that exclusion using an explicitly proved positive initial
coefficient.

\paragraph{Why the determinant-character convention matters.}
The intrinsic grading action is the one detected by Hilbert poles
and by the fixed-locus argument in the plethysm refutation.
It is invariant under the graded-algebra isomorphisms just proved.
A raw $GL$ action on a newly twisted Cartan cone can acquire a
different central kernel, even when its effective graded quotient
is unchanged.  On a preinvariant Cartan tensor product the whole
center can instead act trivially.  These are different stabilizer
questions.  We therefore make no assertion that PNEW-008 bounds
the determinant image in every such raw presentation.  A converse
formulated using that presentation-dependent determinant image
requires an explicit normalization before it has the intrinsic
meaning suggested by Hilbert-series poles.

\paragraph{The broad density clause.}
GCT VI, Hypothesis 3.3.2(b), printed p.~51, concerns specified
subgroup-restriction inputs of length at most $N$.  It allows
ordered products of groups and explicit polynomial maps, with
additive costs as in Sections 3.4.1--3.4.2, printed pp.~53--57.
The inherited P-F63 proof reproduced below, including its
finite-coding assumption, has been reread in full and checked
against those source rules.  It refutes this broad clause under
every finite source-compatible coding that admits constant-cost
adjoining of one fixed ordered factor and has
$\exp(O(N\log(N+2)))$ inputs of source cost at most $N$.
It does not supply an ordinary plethysm padding map, and therefore
does not change the unresolved narrow conclusion above.
